% ============================================================
%  MTH636 – Game Theory
%  Exam Survival Guide / Crash Course
% ============================================================
\documentclass[12pt, a4paper]{article}

% ============================================================
%  MTH636 -- Game Theory
%  Notes Preamble — Color-Coded, TikZ-Enhanced
% ============================================================

% --- Core packages ---
\usepackage[T1]{fontenc}
\usepackage[utf8]{inputenc}
\usepackage{lmodern}
\usepackage[margin=2.5cm, top=2.8cm, bottom=2.8cm]{geometry}
\usepackage{microtype}

% --- Math ---
\usepackage{amsmath, amssymb, amsthm, mathtools}
\usepackage{bm}        % bold math
\usepackage{cancel}    % strike-through in math

% --- Color & Boxes ---
\usepackage[dvipsnames, svgnames, x11names, table]{xcolor}
\usepackage[most]{tcolorbox}
\tcbuselibrary{breakable, skins, theorems}

% --- TikZ & Diagrams ---
\usepackage{tikz}
\usepackage{tikz-qtree}
\usetikzlibrary{arrows.meta, positioning, shapes.geometric, fit, backgrounds,
                decorations.pathreplacing, calc, matrix, trees}
\usepackage{pgfplots}
\pgfplotsset{compat=1.18}

% --- Tables ---
\usepackage{booktabs}
\usepackage{multirow}
\usepackage{array}
\usepackage{tabularx}
\usepackage{makecell}

% --- Typography & Layout ---
\usepackage{enumitem}
\usepackage{parskip}
\usepackage{fancyhdr}
\usepackage{titlesec}
\usepackage{titletoc}
\usepackage{hyperref}
\usepackage{graphicx}
\usepackage{float}
\usepackage{caption}

% --- Color Palette ---
\definecolor{defblue}{HTML}{1A5276}      % Definition header
\definecolor{defbluebg}{HTML}{D6EAF8}    % Definition background
\definecolor{thmgreen}{HTML}{1E8449}     % Theorem header
\definecolor{thmbg}{HTML}{D5F5E3}        % Theorem background
\definecolor{exorange}{HTML}{B7770D}     % Example header
\definecolor{exbg}{HTML}{FEF9E7}         % Example background
\definecolor{warnred}{HTML}{922B21}      % Warning/Key Insight header
\definecolor{warnbg}{HTML}{FADBD8}       % Warning background
\definecolor{noticegray}{HTML}{2C3E50}   % Note header
\definecolor{noticebg}{HTML}{EBF5FB}     % Note background
\definecolor{taketawayyellow}{HTML}{7D6608} % Key Takeaway header
\definecolor{takeawaybg}{HTML}{FDFEFE}   % Takeaway background
\definecolor{accentpurple}{HTML}{6C3483} % Accent / proof
\definecolor{lightgray}{HTML}{F2F3F4}    % shading

% ============================================================
%  TCOLORBOX ENVIRONMENTS
% ============================================================

% -- Definition (Blue) --
\newtcolorbox{defbox}[2][]{
  enhanced, breakable,
  colback=defbluebg, colframe=defblue, coltitle=white,
  fonttitle=\bfseries\sffamily,
  title={Definition: #2},
  attach boxed title to top left={yshift=-2mm, xshift=4mm},
  boxed title style={colback=defblue, rounded corners},
  arc=4pt, boxrule=1.2pt,
  left=6pt, right=6pt, top=8pt, bottom=6pt,
  #1
}

% -- Theorem / Proposition (Green) --
\newtcolorbox{thmbox}[2][]{
  enhanced, breakable,
  colback=thmbg, colframe=thmgreen, coltitle=white,
  fonttitle=\bfseries\sffamily,
  title={Theorem: #2},
  attach boxed title to top left={yshift=-2mm, xshift=4mm},
  boxed title style={colback=thmgreen, rounded corners},
  arc=4pt, boxrule=1.2pt,
  left=6pt, right=6pt, top=8pt, bottom=6pt,
  #1
}

% -- Proposition (Green variant) --
\newtcolorbox{propbox}[2][]{
  enhanced, breakable,
  colback=thmbg, colframe=thmgreen!70!black, coltitle=white,
  fonttitle=\bfseries\sffamily,
  title={Proposition: #2},
  attach boxed title to top left={yshift=-2mm, xshift=4mm},
  boxed title style={colback=thmgreen!70!black, rounded corners},
  arc=4pt, boxrule=1.2pt,
  left=6pt, right=6pt, top=8pt, bottom=6pt,
  #1
}

% -- Example (Orange/Yellow) --
\newtcolorbox{exbox}[2][]{
  enhanced, breakable,
  colback=exbg, colframe=exorange, coltitle=white,
  fonttitle=\bfseries\sffamily,
  title={Example: #2},
  attach boxed title to top left={yshift=-2mm, xshift=4mm},
  boxed title style={colback=exorange, rounded corners},
  arc=4pt, boxrule=1.2pt,
  left=6pt, right=6pt, top=8pt, bottom=6pt,
  #1
}

% -- Key Insight / Warning (Red) --
\newtcolorbox{warnbox}[2][]{
  enhanced, breakable,
  colback=warnbg, colframe=warnred, coltitle=white,
  fonttitle=\bfseries\sffamily,
  title={Key Insight: #2},
  attach boxed title to top left={yshift=-2mm, xshift=4mm},
  boxed title style={colback=warnred, rounded corners},
  arc=4pt, boxrule=1.2pt,
  left=6pt, right=6pt, top=8pt, bottom=6pt,
  #1
}

% -- Note (Gray-Blue) --
\newtcolorbox{notebox}[1][]{
  enhanced, breakable,
  colback=noticebg, colframe=noticegray!60, coltitle=white,
  fonttitle=\bfseries\sffamily,
  title={Note},
  attach boxed title to top left={yshift=-2mm, xshift=4mm},
  boxed title style={colback=noticegray!60, rounded corners},
  arc=4pt, boxrule=0.8pt,
  left=6pt, right=6pt, top=8pt, bottom=6pt,
  #1
}

% -- Key Takeaway --
\newtcolorbox{takebox}[1][]{
  enhanced, breakable,
  colback=takeawaybg, colframe=taketawayyellow, coltitle=white,
  fonttitle=\bfseries\sffamily,
  title={Key Takeaway},
  attach boxed title to top left={yshift=-2mm, xshift=4mm},
  boxed title style={colback=taketawayyellow, rounded corners},
  arc=4pt, boxrule=1.2pt,
  left=6pt, right=6pt, top=8pt, bottom=6pt,
  #1
}

% -- Intuition --
\newtcolorbox{intubox}[1][]{
  enhanced, breakable,
  colback=accentpurple!8, colframe=accentpurple!60, coltitle=white,
  fonttitle=\bfseries\sffamily,
  title={Intuition},
  attach boxed title to top left={yshift=-2mm, xshift=4mm},
  boxed title style={colback=accentpurple!80, rounded corners},
  arc=4pt, boxrule=0.8pt,
  left=6pt, right=6pt, top=8pt, bottom=6pt,
  #1
}

% ============================================================
%  PAYOFF MATRIX HELPER
%  Usage: inside tabular, use \payoff{a}{b} for (a,b)
% ============================================================
\newcommand{\payoff}[2]{$(#1,\;#2)$}
\newcommand{\nash}[2]{\underline{$#1$},\underline{$#2$}}  % underline NE payoffs
\newcommand{\Nset}{\mathbb{N}}
\newcommand{\Rset}{\mathbb{R}}
\newcommand{\Sset}{\mathcal{S}}
\newcommand{\Ui}{u_i}

% ============================================================
%  SECTION FORMATTING
% ============================================================
\titleformat{\section}
  {\large\bfseries\sffamily\color{defblue}}
  {\thesection}{1em}{}[\titlerule]

\titleformat{\subsection}
  {\normalsize\bfseries\sffamily\color{thmgreen!80!black}}
  {\thesubsection}{1em}{}

\titleformat{\subsubsection}
  {\normalsize\itshape\sffamily\color{exorange!80!black}}
  {\thesubsubsection}{1em}{}

% ============================================================
%  HEADER / FOOTER
% ============================================================
\pagestyle{fancy}
\fancyhf{}
\fancyhead[L]{\sffamily\small\color{defblue} MTH636 -- Game Theory}
\fancyhead[R]{\sffamily\small\color{defblue} \leftmark}
\fancyfoot[C]{\sffamily\small\thepage}
\renewcommand{\headrulewidth}{0.4pt}

% ============================================================
%  HYPERREF SETUP
% ============================================================
\hypersetup{
  colorlinks=true,
  linkcolor=defblue,
  urlcolor=thmgreen,
  citecolor=accentpurple,
  bookmarksnumbered=true,
  pdftitle={MTH636 Game Theory Notes},
  pdfauthor={MTH636 Student Notes}
}

% ============================================================
%  MISC HELPERS
% ============================================================
\newcommand{\BR}{\text{BR}}        % Best Response
\newcommand{\NE}{\text{NE}}        % Nash Equilibrium
\newcommand{\SPNE}{\text{SPNE}}    % Subgame Perfect NE
\newcommand{\Dom}{\text{dom}}      % Dominance
\newcommand{\mixs}[1]{\sigma_{#1}} % Mixed strategy
\newcommand{\pures}[1]{s_{#1}}     % Pure strategy


% ============================================================
\begin{document}

% Title
\begin{center}
  {\LARGE\bfseries\sffamily\color{defblue} Game Theory: The Exam Survival Guide}\\[6pt]
  {\Large\sffamily (How to Solve Exam Questions if You Skipped the Lectures)}\\[4pt]
\end{center}
\vspace{0.3cm}
\hrule height 1.5pt
\vspace{0.6cm}

This guide is designed to teach you the mechanical skills required to solve Game Theory exam questions from scratch. If you skipped a month of lectures, read this carefully. 

% ============================================================
\section{Module 1: Simultaneous Games (Matrices)}

You will be given a grid (a matrix). \textbf{Player 1} chooses the \textbf{Row} (Top, Bottom). \textbf{Player 2} chooses the \textbf{Column} (Left, Right). The numbers in the cells are payoffs: \texttt{(Player 1's payoff, Player 2's payoff)}.

\subsection{Skill 1: Eliminating Dominated Strategies (IESDS)}
Before solving anything, simplify the game. If a player has a move that is \textit{always} worse than another move, they will never play it.

\textbf{The Algorithm:}
\begin{enumerate}
    \item \textbf{Check P1 (Rows):} Look at Row A and Row B. If EVERY number for P1 in Row A is strictly less than the number directly below it in Row B, cross out Row A forever. (Note: A mixed strategy, like 50\% Row B and 50\% Row C, can also dominate Row A!).
    \item \textbf{Check P2 (Columns):} Look at Col X and Col Y. If EVERY number for P2 in Col X is strictly less than the number directly next to it in Col Y, cross out Col X forever.
    \item Repeat until you can't cross out anything else.
\end{enumerate}

\subsection{Skill 2: Finding Pure Strategy Nash Equilibria}
A Nash Equilibrium (NE) is a cell where neither player regrets their choice after seeing the other's choice.

\textbf{The Underline Algorithm (Foolproof):}
\begin{enumerate}
    \item \textbf{Pretend you are Player 1.} Look at the \textit{first column}. Which cell gives Player 1 the highest number? Underline that number. Repeat for every column.
    \item \textbf{Pretend you are Player 2.} Look at the \textit{first row}. Which cell gives Player 2 the highest number? Underline that number. Repeat for every row.
    \item Look at the matrix. Any cell that has \textbf{BOTH numbers underlined} is a Pure Strategy Nash Equilibrium. Report it as \texttt{(Row Name, Column Name)}.
\end{enumerate}

% ============================================================
\section{Module 2: Mixed Strategy Nash Equilibria}

\textbf{When to use this:} If the underline method gives you zero cells with both numbers underlined (like Rock-Paper-Scissors), or if the exam explicitly asks for \textit{all} Nash Equilibria (including mixed ones).
\textbf{The Core Concept:} You randomize your moves using probabilities to make your opponent completely indifferent to what they choose. (Nash's Theorem guarantees at least one NE always exists in finite games).

\subsection{Skill 3: The 2x2 Mixed Strategy Algorithm}
Let's solve a generic $2 \times 2$ game.
\begin{center}
\renewcommand{\arraystretch}{1.4}
\begin{tabular}{c|cc}
  $P_1 \setminus P_2$ & L (prob $q$) & R (prob $1-q$) \\
  \hline
  T (prob $p$) & \payoff{a}{b} & \payoff{c}{d} \\
  B (prob $1-p$) & \payoff{e}{f} & \payoff{g}{h} \\
\end{tabular}
\end{center}

\textbf{Step 1: Find P2's probabilities ($q$ and $1-q$)} \\
We must make Player 1 indifferent between playing Top and Bottom.
\begin{itemize}
    \item If P1 plays Top, their expected payoff is: $E[T] = a(q) + c(1-q)$
    \item If P1 plays Bottom, their expected payoff is: $E[B] = e(q) + g(1-q)$
    \item \textbf{Set them equal:} $a(q) + c(1-q) = e(q) + g(1-q)$.
    \item Solve for $q$ using basic algebra. If $q = 1/3$, then P2 plays L 33\% of the time and R 67\% of the time.
\end{itemize}

\textbf{Step 2: Find P1's probabilities ($p$ and $1-p$)} \\
We must make Player 2 indifferent between playing Left and Right.
\begin{itemize}
    \item If P2 plays Left, their expected payoff is: $E[L] = b(p) + f(1-p)$
    \item If P2 plays Right, their expected payoff is: $E[R] = d(p) + h(1-p)$
    \item \textbf{Set them equal:} $b(p) + f(1-p) = d(p) + h(1-p)$.
    \item Solve for $p$.
\end{itemize}

\textbf{Step 3: State the Equilibrium.} \\
The mixed NE is $\sigma_1 = (p, 1-p)$ and $\sigma_2 = (q, 1-q)$. \\
\textit{Exam Check:} If your $p$ or $q$ comes out negative or greater than 1, you made an algebra mistake, or that specific mix isn't an equilibrium!

% ============================================================
\section{Module 3: Sequential Games (Extensive Form)}

You will be given a \textbf{Game Tree}. Players take turns making decisions at ``nodes''.
\begin{itemize}
    \item \textbf{Perfect Information:} You always know exactly where you are in the tree.
    \item \textbf{Imperfect Information (Information Sets):} A dotted bubble encloses two or more nodes. It means the player whose turn it is \textit{cannot distinguish} which node they are at (e.g., they don't know what the previous player did). Treat the bubble like a Simultaneous Matrix game!
\end{itemize}

\subsection{Skill 4: Backward Induction (Finding SPNE)}
A \textbf{Subgame Perfect Nash Equilibrium (SPNE)} is a sequence of moves where no one is making a ``non-credible threat'' (an empty promise they wouldn't actually want to carry out if forced to).

\textbf{The Algorithm (Folding the Tree):}
\begin{enumerate}
    \item Start at the very bottom of the tree (the final choices before the payoffs).
    \item Look at the player making the choice at that node. Which branch gives \textbf{them} the higher payoff?
    \item \textbf{Highlight that branch.} Cross out the other branch.
    \item Mentally ``delete'' that decision node and replace it entirely with the payoff of the highlighted branch.
    \item Move up one level to the previous player's decision. Repeat the process, choosing the best payoff for \textit{that} player based on the new collapsed tree.
    \item Continue until you reach the root (start) of the tree. The path of highlighted branches is your SPNE!
\end{enumerate}

% ============================================================
\section{Module 4: Zero-Sum Games (Minimax \& Saddle Points)}

In a zero-sum game, whatever P1 wins, P2 loses. The matrix usually only shows P1's payoffs (because P2's payoff is just the negative of P1's). P1 wants to MAXIMIZE the number, P2 wants to MINIMIZE the number.

\subsection{Skill 5: Finding a Saddle Point (Pure Value)}
\begin{enumerate}
    \item \textbf{Row Mins:} For each row, find the \textit{minimum} value. Write it next to the row. P1 will choose the \textbf{MAXIMUM of these row minimums} (Max-Min).
    \item \textbf{Col Maxes:} For each column, find the \textit{maximum} value. Write it under the column. P2 will choose the \textbf{MINIMUM of these column maximums} (Min-Max).
    \item If the Max-Min equals the Min-Max, that number is the \textbf{Value of the Game}, and the cell where they intersect is a \textbf{Saddle Point} (which is a Pure NE). 
    \item If they do NOT equal, the game has no pure value. You must use the Mixed Strategy algorithm (Skill 3) to find the mixed value!
\end{enumerate}

% ============================================================
\section{Module 5: Continuous Strategy Games (Calculus)}

Sometimes players don't choose from a matrix; they choose a number (e.g., how much effort to put in, what price to set, or where to stand on a line from 0 to 1). 

\subsection{Skill 6: Solving with Derivatives (Best Response Functions)}
If an exam question gives you a utility equation like $U_1(e_1, e_2) = e_1 e_2 - c e_1$, do NOT panic. 

\textbf{The Algorithm:}
\begin{enumerate}
    \item \textbf{Write P1's Utility:} $U_1$.
    \item \textbf{Take the partial derivative} with respect to P1's choice ($e_1$): $\frac{\partial U_1}{\partial e_1}$.
    \item \textbf{Set it to 0 and solve for $e_1$:} This gives you Player 1's \textbf{Best Response Function}, $e_1^*(e_2)$. It tells P1 exactly what to do for any choice P2 makes.
    \item \textbf{Repeat for P2:} Take $\frac{\partial U_2}{\partial e_2} = 0$ and solve for $e_2^*(e_1)$.
    \item \textbf{Solve the System:} You now have two equations with two variables. Substitute one into the other to find the exact numbers. That pair $(e_1^*, e_2^*)$ is your Nash Equilibrium!
\end{enumerate}

\subsection{Skill 7: N-Player Symmetric Games}
If there are $n$ players and they are all identical (like $n$ people witnessing a crime):
\begin{itemize}
    \item A \textbf{symmetric Pure NE} is when everyone does the exact same thing (Everyone calls the police, or Nobody calls). Just plug in those scenarios and check if 1 person would want to deviate.
    \item A \textbf{symmetric Mixed NE} is when everyone plays the same probability $p$. Set the expected utility of Action A equal to Action B for Player 1, assuming all other $(n-1)$ players play $p$. Use the Binomial distribution for probabilities!
\end{itemize}

\end{document}
